\section{An explicit formula for height five}\label{sec:five}

At height five the transfer calculation reduces to a
six-dimensional polynomial certificate.
We give the reduced matrix and boundary vectors explicitly;
Appendix \ref{app:certificate} supplies the single-column
matrices from which they are obtained.

Represent a row subset by a five-bit mask, with the least
significant bit corresponding to the top row.
The row-reversal orbits of reachable even-column states,
in the chosen order, are
\begin{equation}\label{eq:even-orbits}
 \{0\},\ \{3,24\},\ \{6,12\},\ \{9,18\},\
 \{15,30\},\ \{27\}.
\end{equation}
Row-reversed states have equal counts.
Let each component of $u_j$ record the count of \emph{one
individual state} in the corresponding orbit after $2j$
columns, rather than the sum over that orbit.
Then $u_0=\evec$ and $u_{j+1}=Mu_j$, where
\begin{equation}\label{eq:M5}
 M=\begin{pmatrix}
8&6&4&2&2&1\\
3&4&1&1&1&1\\
2&1&2&1&1&0\\
1&1&1&2&1&0\\
1&1&1&1&1&0\\
1&2&0&0&0&1
\end{pmatrix}.
\end{equation}
Define
\begin{equation}\label{eq:PN}
 P(t)=1-15t+32t^2-15t^3+t^4,\qquad
 N(t)=\begin{pmatrix}
1-7t+7t^2-t^3\\
3t-4t^2+t^3\\
2t-5t^2+t^3\\
t+t^2\\
t-t^3\\
t-t^3
\end{pmatrix}.
\end{equation}
Polynomial multiplication gives the exact identity
\begin{equation}\label{eq:certificate}
 (I-tM)N(t)=P(t)\evec.
\end{equation}
Since $P(0)=1$, it follows that
\[
 f(t):=\sum_{j\ge0}u_jt^j=\frac{N(t)}{P(t)}.
\]
Appendix \ref{app:certificate} also records all four nonconstant
coefficient-vector identities underlying
\eqref{eq:certificate}.

The terminal row vectors for vertical-axis reflection are
\[
 w_{v,e}=(1,2,2,2,2,1),\qquad
 w_{v,o}=(8,6,4,4,2,1),
\]
and those for the half-turn are
\[
 w_{r,e}=(1,0,0,0,0,1),\qquad
 w_{r,o}=(2,2,0,0,0,1).
\]
The even vectors sum all states or only row-reversal-fixed
states. The odd vectors first apply one column transition
and then take the appropriate sum.
Lemma \ref{lem:half} gives
\begin{equation}\label{eq:terminal5}
\begin{split}
 V_5(x)&=(w_{v,e}+xw_{v,o})f(x^2),\\
 R_5(x)&=(w_{r,e}+xw_{r,o})f(x^2).
\end{split}
\end{equation}

\begin{theorem}\label{thm:five}
The four fixed-point generating functions at height five are
\begin{align}
 D_5(x)&=\frac{1-7x+7x^2-x^3}{P(x)},&
 H_5(x)&=\frac{1-x}{1-3x+x^2},\label{eq:five-DH}\\
 V_5(x)&=\frac{1+8x+8x^2-23x^3-9x^4+16x^5-x^7}{P(x^2)},\label{eq:five-V}\\
 R_5(x)&=\frac{1+2x-6x^2-7x^3+7x^4+6x^5-2x^6-x^7}{P(x^2)}.\label{eq:five-R}
\end{align}
The orbit generating function is
\begin{equation}\label{eq:five-A}
 A_5(x)=\frac{D_5(x)+H_5(x)+V_5(x)+R_5(x)}4.
\end{equation}
\end{theorem}
\begin{proof}
A complete tiling has empty terminal state, so $D_5=f_0$,
giving the first formula.
Substitution of \eqref{eq:PN} into \eqref{eq:terminal5} gives
the two formulas for $V_5,R_5$.

By Lemma \ref{lem:h-odd}, $h_{5,n}=t_2(2n)$.
A two-row tiling either starts with one vertical domino or
with two horizontal dominoes covering its first two columns.
Thus $t_2(j)=t_2(j-1)+t_2(j-2)$ for $j\ge2$, with initial
values $t_2(0)=t_2(1)=1$.
The two-step recurrence matrix is
$L=\left(\begin{smallmatrix}2&1\\1&1\end{smallmatrix}\right)$,
which satisfies $L^2-3L+I=0$.
Hence $c_n=t_2(2n)$ satisfies $c_{n+2}=3c_{n+1}-c_n$,
with $c_0=1,c_1=2$, proving the formula for $H_5$.
Finally, Burnside's formula proves \eqref{eq:five-A}.
\end{proof}

Write
\[
 Q(x)=(1-3x+x^2)P(x)P(x^2)=\sum_{j=0}^{14}q_jx^j,\qquad
 U(x)=\sum_{j=0}^{13}b_jx^j,
\]
where $A_5(x)=U(x)/Q(x)$.
Table \ref{tab:five-coeff} lists all coefficients and the
fourteen initial values needed for
\begin{equation}\label{eq:rec5}
 a_5(n)=-\sum_{j=1}^{14}q_j a_5(n-j),\qquad n\ge14.
\end{equation}
No minimality assertion about the recurrence order is needed.

\begin{table}[htbp]
\centering
\caption{Coefficients of $U,Q$ and initial values; $b_{14}=0$.}\label{tab:five-coeff}
\begin{tabular}{rrrr}
\toprule
$j$ & $b_j$ & $q_j$ & $a_5(j)$\\
\midrule
0 & 1 & 1 & 1\\
1 & -13 & -18 & 5\\
2 & 6 & 63 & 33\\
3 & 194 & 144 & 329\\
4 & -365 & -1060 & 3818\\
5 & -371 & 1296 & 46878\\
6 & 1304 & 1312 & 584386\\
7 & -594 & -3492 & 7318152\\
8 & -703 & 1312 & 91752831\\
9 & 717 & 1296 & 1150774799\\
10 & -154 & -1060 & 14434490195\\
11 & -36 & 144 & 181060839597\\
12 & 15 & 63 & 2271175917429\\
13 & -1 & -18 & 28489049801109\\
14 & 0 & 1 & \text{by recurrence}\\
\bottomrule
\end{tabular}
\end{table}


The values $a_5(1),\ldots,a_5(5)=5,33,329,3818,46878$
agree with the entries for widths $2,4,6,8,10$ in
\cite[Table 12]{mathar2013}.
This is a comparison made after the derivation; those values
were not used to fit a matrix or a generating function.
